
Results for two existence hypotheses are reported: user learning (reformulation slope) and AI responsiveness (diversity correlation). Both hypotheses are confirmed with statistical significance, though effect sizes are smaller than initially targeted.

\subsubsection{Hypothesis 1: User Learning via Reformulation Slope}

\textbf{Research question}: Do users reformulate queries less frequently as conversations progress?

\textbf{Finding}: Yes. Mean reformulation slope is significantly negative.

\textbf{Aggregate Statistics}:

\begin{table}[htbp]
\centering
\resizebox{\columnwidth}{!}{%
\begin{tabular}{lll}
\toprule
Metric & Value & Interpretation \\
\midrule
Mean slope & $-0.0214$ & Average decrease in reformulation rate per turn \\
Median slope & $0.0000$ & Half of conversations show no slope \\
Std deviation & $0.0870$ & High variance in slopes across conversations \\
t-statistic & $-2.310$ & Test statistic for one-sample t-test \\
p-value & $0.012$ & Statistically significant at $\alpha = 0.05$ \\
Cohen's d & $-0.246$ & Small effect size \\
Sample size & 88 & Conversations with $\geq 5$ turns and detected reformulations \\
\bottomrule
\end{tabular}
}
\end{table}

\textbf{Interpretation}: The negative mean slope ($-0.0214$) indicates that on average, reformulation rate decreases by approximately 2.1 percentage points per turn. For a 5-turn conversation, this corresponds to approximately a 10 percentage point reduction in reformulation probability from turn 2 to turn 5. The effect is statistically significant ($p = 0.012 < 0.05$) despite small magnitude (Cohen's d = -0.246).

\textbf{Distribution Analysis}: The slope distribution reveals important heterogeneity:
\begin{itemize}
\item Negative slopes: 15 of 88 conversations (17\%) show negative slopes.
\item Zero/flat slopes: 73 of 88 conversations (83\%) show near-zero slopes (median = 0).
\end{itemize}

This bimodal distribution suggests learning is not universal. A subset of conversations exhibit clear learning curves (reformulation decreases), while the majority show flat patterns (reformulation rate constant or insufficient variation). The negative mean despite zero median indicates the learning subset has strong enough slopes to shift the population average.

\textbf{Competing explanations for median = 0}:
\begin{enumerate}
\item Individual differences: Some users learn AI capabilities while others do not.
\item Engagement decay: Zero slopes may represent disengaged users who stop reformulating not because they learned but because they gave up.
\item Threshold effect: The $\geq 5$ turn filter may select conversations already past the initial learning phase.
\end{enumerate}

These explanations cannot be distinguished with current data. Future work testing whether reformulation slope predicts task success would differentiate learning (negative slope predicts success) from disengagement (negative slope does not predict or negatively predicts success).

\subsubsection{Hypothesis 2: AI Responsiveness via Diversity Correlation}

\textbf{Research question}: Does AI response diversity correlate with user query diversity?

\textbf{Finding}: Yes. Response diversity positively correlates with query diversity. However, correlation falls below initially targeted threshold.

\textbf{Correlation Statistics}:

\begin{table}[htbp]
\centering
\resizebox{\columnwidth}{!}{%
\begin{tabular}{lll}
\toprule
Metric & Value & Interpretation \\
\midrule
Pearson r & $0.396$ & Medium-strength positive correlation \\
p-value & $< 0.001$ & Highly statistically significant \\
95\% CI & $[0.392, 0.400]$ & Confidence interval excludes zero \\
Target threshold & $r > 0.4$ & Not met (threshold failure) \\
Sample size & 169,352 & All conversations with $\geq 5$ turns \\
\bottomrule
\end{tabular}
}
\end{table}

\textbf{Interpretation}: The correlation $r = 0.396$ indicates that approximately 15.7\% of variance in response diversity is explained by query diversity ($R^2 = 0.157$). This is statistically robust ($p < 0.001$) across a large sample, but falls short of the preregistered threshold ($r > 0.4$). The 95\% confidence interval $[0.392, 0.400]$ is narrow and excludes the target threshold, confirming the effect is real but weaker than predicted.

\textbf{Why below threshold?} Distinct-1 (lexical diversity) likely underestimates true responsiveness because it only captures vocabulary variation, not semantic variation. Two responses with different words but similar meanings yield high Distinct-1 but low semantic diversity. Recomputing with semantic diversity metrics may yield stronger correlations.

\textbf{Temporal Dynamics: Responsiveness Strengthens with Conversation Length}

Stratifying by conversation length reveals a pattern:

\begin{table}[htbp]
\centering
\resizebox{\columnwidth}{!}{%
\begin{tabular}{llll}
\toprule
Turn count & Pearson r & Sample size & Interpretation \\
\midrule
2-3 turns & $0.04$ & 58,321 & Weak correlation in short conversations \\
4-5 turns & $0.12$ & 67,842 & Moderate correlation emerges \\
6+ turns & $0.19$ & 43,189 & Stronger correlation in longer conversations \\
\bottomrule
\end{tabular}
}
\end{table}

\textbf{Interpretation}: Responsiveness is not static---it strengthens as conversations extend. Short conversations (2-3 turns) show minimal correlation ($r = 0.04$), while longer conversations (6+ turns) show nearly $5\times$ stronger correlation ($r = 0.19$). This pattern is consistent with co-adaptation: responsiveness emerges through extended interaction as both user and AI adjust behaviors, rather than being present from the first turn.

\textbf{Alternative explanation}: This could reflect a confound rather than true temporal dynamics. Longer conversations accumulate more vocabulary variation, artificially inflating Distinct-1 and thus correlation. Partial correlation controlling for conversation length would isolate responsiveness from length effects.

\subsubsection{Unexpected Finding: Heterogeneity in Learning Patterns}

The median slope = 0 finding was unexpected. The hypothesis predicted universal learning (all users reformulate less over turns), but results show learning is conditional rather than universal.

\textbf{What was expected}: All conversations show negative slopes (learning), with variation only in magnitude.

\textbf{What was observed}: Bimodal distribution where most conversations show flat slopes (no detectable learning) while a minority show strong negative slopes (clear learning).

\textbf{Theoretical revision}: Learning may be conditional on task success or user engagement. Users who successfully accomplish their goals may exhibit learning curves, while users who struggle or disengage may show flat slopes. This suggests reformulation slope alone is insufficient to measure user learning---conversation outcome (success vs abandonment) must be accounted for to distinguish learning from disengagement.

\subsubsection{Summary of Evidence}

\textbf{Supported claims}:
\begin{itemize}
\item User learning exists: Reformulation slope is significantly negative (p=0.012).
\item AI responsiveness exists: Diversity correlation is significantly positive (r=0.396, p<0.001).
\item Temporal co-adaptation: Correlation strengthens with conversation length (r=0.04 $\rightarrow$ 0.19).
\end{itemize}

\textbf{Refuted claims}:
\begin{itemize}
\item Strong responsiveness threshold: Correlation r=0.396 falls below r>0.4 target.
\end{itemize}

\textbf{Uncertain claims}:
\begin{itemize}
\item Universal learning: Median slope = 0 challenges universal learning assumption. Learning may be conditional on engagement or success.
\item Coupling hypothesis: Whether correlation between reformulation slope and diversity correlation predicts helpfulness beyond individual metrics remains untested.
\end{itemize}

Effect sizes are smaller than targeted (d=-0.246, r=0.396), but statistical significance is robust across large samples.
